\documentclass[11pt]{article}
% Shared preamble for per-Python-file documentation (input via \input{...}).
\usepackage[margin=1in]{geometry}
\usepackage[T1]{fontenc}
\usepackage[utf8]{inputenc}
\usepackage{lmodern}
\usepackage{hyperref}
\usepackage{xcolor}
\usepackage{listings}
\usepackage{listingsutf8}
\lstset{
  basicstyle=\ttfamily\small,
  columns=fullflexible,
  breaklines=true,
  upquote=true,
  inputencoding=utf8,
  extendedchars=true,
  frame=single,
  framerule=0.2pt,
  tabsize=2
}


\title{\texttt{utils.py} (Q\&A)}
\author{}
\date{}

\begin{document}
\maketitle

\paragraph{Q: What is the purpose of \texttt{utils.py}?}
\textbf{A:} This file provides shared evaluation metrics and helper routines used across multiple run scripts (notably MAE, XAUC, and KL divergence), plus utilities used by the TPM baseline for tree-style label encoding and decoding.

\paragraph{Q: What are the core evaluation metrics implemented here?}
\textbf{A:} The main metrics are \lstinline|eval_mae| (mean absolute error), \lstinline|eval_xauc| (a ranking-style agreement metric based on ordering induced by predictions), and \lstinline|eval_kl| (KL divergence between binned empirical distributions of labels vs. predictions).

\paragraph{Q: How does XAUC work at a high level in this implementation?}
\textbf{A:} The code sorts (label, prediction) pairs by prediction descending, then counts how many label inversions occur in that sorted list (using \lstinline|InversePairsCalc|). XAUC is computed as inversions divided by the total number of pairs, effectively measuring whether larger predictions correspond to larger labels.

\paragraph{Q: What TPM-specific utilities are implemented here?}
\textbf{A:} The TPM pipeline requires (1) building percentile bucket boundaries (\lstinline|get_playtime_percentiles_range|), (2) encoding labels into a binary tree of interval decisions (\lstinline|get_tree_encoded_label|), (3) decoding model outputs back to a scalar value and an uncertainty-like term (\lstinline|get_tree_encoded_value|), and (4) computing a weighted tree classification loss (\lstinline|get_tree_classify_loss|).

\paragraph{Q: What are the key code entrypoints a reader should look at first?}
\textbf{A:} If you are focused on evaluation: \lstinline|eval_mae|, \lstinline|eval_xauc|, and \lstinline|eval_kl|. If you are focused on TPM: \lstinline|get_playtime_percentiles_range|, \lstinline|get_tree_encoded_label|, \lstinline|get_tree_encoded_value|, and \lstinline|get_tree_classify_loss|.

\paragraph{Q: Code example: how a run script typically uses these metrics?}
\textbf{A:}
\begin{lstlisting}[language=Python]
mae = eval_mae(labels, scores)
xauc = eval_xauc(labels, scores)
kl = eval_kl(labels, scores)
\end{lstlisting}

\paragraph{Q: Full source code?}
\textbf{A:}
\lstinputlisting[language=Python]{/workspace/utils.py}

\end{document}
